\section{Discussion}

The baseline reproduces the core optimization mechanics: sampled graphs remain acyclic, edge logits receive finite gradients, and truthful-to-decision probabilities increase above 0.63 for every seed. It does not reproduce an optimized-over-random accuracy gain. The equal 0.250 means and unstable per-seed difference make the honest conclusion a null average result under this reduced setup.

Full and optimized graphs outperform direct answer on average, but full-graph variance comes from seeded majority-vote ties between truthful and adversarial answers. Random graphs use fewer tokens partly because disconnection prevents model calls and valid decisions. Resource comparisons must therefore report failures jointly with tokens.

The baseline result strengthens the motivation for explicit feasibility in BC-GPTSwarm: a method should not satisfy a resource target by disconnecting the decision node. Phase 5 addresses this mechanically by conditioning the policy on a valid decision input and correcting the conditional log probability. The mock smoke tier shows that the implementation can reduce sampled cost without invalid graphs, but it does not establish real-model budget control or accuracy retention. The later comparison must include fixed-penalty and accuracy-only controls rather than assuming that resource pressure improves reasoning. \textbf{TODO: Update after proposed-method experiments with budget compliance, trade-offs, and failure cases.}
